\section{Related Work}
\label{sec:related}

\subsection{Self-Supervised Learning for Time Series}
\label{sec:related_ssl}

Self-supervised pretraining for time series has followed two paradigms: contrastive and generative. Both have matured rapidly, yet neither has been adapted for mechanical prognostics.

\paragraph{Contrastive methods.} \citet{yue2022ts2vec} learn representations via hierarchical contrastive learning across temporal and instance dimensions. \citet{tonekaboni2021tnc} exploit the stationarity of temporal neighborhoods, defining positive pairs within stationary windows---an assumption violated by bearing degradation, which is fundamentally non-stationary. \citet{eldele2021tstcc} combine temporal and contextual contrasting, while \citet{zhang2022tfc} align time-domain and frequency-domain representations, demonstrating cross-domain transfer on medical and mechanical datasets. Of these, TF-C is closest to our setting, but evaluates on fault \emph{classification}, not RUL regression.

\paragraph{Masked reconstruction.} \citet{nie2023patchtst} adapt the Vision Transformer to time series via channel-independent patching and masked pretraining, achieving state-of-the-art long-term forecasting. \citet{dong2023simmtm} reconstruct masked time points from multiple masked versions to recover manifold structure, and \citet{cheng2023timemae} combine masked autoencoders with decoupled representation learning. These methods optimize for reconstruction fidelity, which biases representations toward low-level waveform statistics rather than degradation-relevant spectral features.

\paragraph{Foundation models.} \citet{das2024timesfm}, \citet{ansari2024chronos}, \citet{goswami2024moment}, and \citet{gao2024units} pretrain large models on diverse time series corpora for zero-shot or multi-task forecasting. These foundation models predict the \emph{next value}---a fundamentally different objective from encoding degradation state. Our work targets this gap: rather than general-purpose time series modeling, we learn representations aligned with mechanical degradation indicators such as spectral centroid shift.

\subsection{Remaining Useful Life Prediction}
\label{sec:related_rul}

RUL prediction has progressed from physics-based models to deep learning. Early work applied CNNs~\citep{li2018dcnn} and LSTMs~\citep{zheng2017lstm} to sensor data; modern architectures combine these with attention mechanisms. On C-MAPSS~\citep{saxena2008cmapss}, \citet{ttsnet2025} achieve the current state of the art by fusing Transformer, TCN, and self-attention branches (FD001 RMSE: 11.02, Score: 194.6).

For bearing prognostics, FEMTO/PRONOSTIA~\citep{nectoux2012pronostia} and XJTU-SY~\citep{wang2018xjtusybearing} are the primary benchmarks. Multi-scale separable convolution Transformers achieve RMSE 0.124 on FEMTO and 0.160 on XJTU-SY~\citep{mdsct2024}. \citet{nomi2024} exploit time-invariant Neural ODE dynamics for cross-domain bearing transfer. However, nearly all methods share two limitations: (1)~they require full temporal trajectory context (ordered sequences of windows), and (2)~they are fully supervised or use domain adaptation with labeled target data. \citet{shen2026dcssl} present DCSSL, the only published SSL method for bearing RUL on FEMTO, using dual-dimensional contrastive pretraining with trajectory context. Our work differs in three respects: we use JEPA (not contrastive) pretraining, operate on single windows without trajectory context, and demonstrate cross-dataset transfer without explicit domain adaptation.

\subsection{Joint Embedding Predictive Architectures}
\label{sec:related_jepa}

Joint Embedding Predictive Architectures (JEPA) learn by predicting representations of masked inputs from visible context in a shared latent space~\citep{lecun2022path}. \citet{assran2023ijepa} demonstrate with I-JEPA that predicting in representation space avoids the low-level reconstruction bias of masked autoencoders. \citet{bardes2024vjepa} extend to video with 90\% spatiotemporal masking, learning motion and dynamics representations. \citet{wang2024brainjepa} apply JEPA to fMRI brain dynamics with domain-specific positioning and masking (NeurIPS 2024 Spotlight), the closest analogue to our work in adapting JEPA to a non-vision domain with physics-informed structure.

For time series, \citet{tsjepa2024} adapt JEPA with temporal masking for classification benchmarks, while \citet{mtsjepa2026} introduce codebook regularization for multivariate series, preventing representation collapse through discrete bottlenecks. \citet{balestriero2025lejepa} replace heuristic collapse prevention (EMA, stop-gradient) with the theoretically grounded SIGReg regularizer, proving that the optimal embedding distribution is isotropic Gaussian. None of these time series JEPA variants evaluate on prognostics or RUL tasks.

Two recent works apply self-supervised pretraining to vibration data. \citet{wang2025rmgpt} use GPT-style generative pretraining on rotating machinery signals (68.5M parameters), achieving 92\% one-shot fault classification but predicting raw signal tokens rather than latent representations. \citet{openMAE2025} apply masked autoencoders to 5M vibration samples from seismic sensors, showing that large-scale vibration pretraining transfers to downstream sensing tasks. However, \emph{neither applies JEPA to mechanical systems, and neither evaluates on RUL prediction}. Our work is the first to evaluate JEPA for mechanical RUL prediction, diagnose why standard JEPA representations have low correlation with degradation indicators (spectral centroid shift), and demonstrate that hybrid architectures overcome this limitation.

\subsection{Learning with Limited Failure Data}
\label{sec:related_rare}

The scarcity of failure data is a central challenge in prognostics~\citep{weak_supervision_phm2025}. Grey swan events---rare but physically plausible failures---are well-studied in weather prediction, where AI models consistently fail on out-of-distribution extremes~\citep{greyswan_weather2024}, and in supply chain risk, where formal frameworks distinguish grey swans (structurally predictable) from black swans (epistemically unknowable)~\citep{greyswan_supplychain2025}. \citet{greyswan_factory2026} generate synthetic grey swans via gradient-based optimization through differentiable physics models. In mechanical prognostics, \citet{physics_augment_rul2024} demonstrate competitive RUL prediction with zero real failure data using physics-informed augmentation, while \citet{deng2025gan_pinn} combine GANs with inverse PINNs to generate physically valid rare fault signals. \citet{digitaltwin_bearing2025} achieve $>94\%$ fault classification with only 5 labeled real samples by pretraining on physics-simulated data. Reviews of physics-informed RUL methods~\citep{physics_rul_review2024} consistently find that physics constraints provide the largest benefit in low-data regimes---precisely the grey swan scenario. Our self-supervised pretraining approach is complementary: by learning general vibration representations from abundant unlabeled data (including healthy operation), we reduce the labeled data requirement to 18 failure trajectories, without requiring a physics simulator or handcrafted augmentation strategies.
